\chapter{Scarcity, access and unequal exposure}\label{ch:access}
\question{Why can needs remain unmet even when some relevant resources exist?}

A dry well, a broken pipe and an unpaid connection fee can all prevent water from reaching a home. They are not the same obstacle. More money cannot instantly refill the well. More rainfall cannot by itself repair the pipe. A working pipe does not settle who is allowed to use it or who can afford the service.

FAO distinguishes physical scarcity, access scarcity and infrastructure scarcity in its account of water. The distinction keeps the availability of suitable water separate from the arrangements needed to deliver it.\cite{water} We can extend the question cautiously in our clinic example: what exactly prevents a specified service from reaching a specified person at a specified time?

\section{Finding the constraint}
Physical availability includes quantity, location, timing and quality. A region may have water in one season and inadequate storage for another. Water can be present but unsuitable for a particular use. In our example these are questions for physical measurement and domain knowledge, not judgments inferred from the patient's bank balance.

Infrastructure includes treatment, pipes, pumps, roads, power and the ability to maintain them. An institution can own abundant equipment that it cannot operate because a small component is missing. Distribution concerns where usable resources go and under what rules. Affordability concerns whether the required payment fits within a person's resources and competing essential expenses.

These constraints may reinforce one another. A broken network may force people to purchase a more expensive substitute. An unaffordable service may remain unused while installed capacity stands idle. Neither observation would prove that physical limits are unreal. It would tell us to identify the active constraint before choosing a remedy.

Food provides the same conceptual distinction without requiring a disputed global count. Food on a warehouse shelf is physical availability at that place. A household's ability to obtain, store, prepare and safely consume it involves further conditions. A delivery record and a purchase record each answer part of that story. Neither is a complete measure of nourishment.

\section{Socioeconomic conditions are several things}
Socioeconomic status, often shortened to SES, describes a person's or household's position through measures such as income, education and occupation. Wealth, housing security, social exclusion and health add related but distinct information. An income number alone does not describe all of these conditions.

The WHO's 2025 report on social determinants of health equity describes how living and working conditions and access to power and resources influence avoidable health gaps.\cite{health} That supports attention to mechanisms. It does not establish that every health difference has a single monetary cause.

Cause and association must also be distinguished. Poor health can reduce earning opportunities; insecure work can make it harder to maintain health; earlier conditions can influence both. A comparison between groups does not automatically tell us what caused one individual's circumstances. To identify an effect, a study needs a design that can address competing explanations.

One bounded example is the US Moving to Opportunity experiment. Families in five cities were randomly offered different housing-assistance options during 1994--1998. A later study found improved adult outcomes for children who moved young, while effects for adolescents differed. This is evidence about a particular intervention and population, not a universal coefficient linking money to opportunity or a test of EBU.\cite{mto}

\section{Two households in the same heatwave}
Imagine two households during the same hot week. One has shade, reliable electricity, a room that stays cooler and flexible work. The other lives under an uninsulated roof, has a member who needs assistance and cannot easily leave work or travel. This is a schematic example of unequal exposure and ability to respond, not a clinical risk calculation.

The same outdoor temperature does not make their situations identical. A physical record limited to outdoor temperature would omit the housing and access differences. A record limited to income would omit other relevant conditions. Expanding a model can make it more informative, but also creates demands for evidence, privacy protection and clear interpretation.

The example illustrates a difficulty for any scalar account. If a person's disability makes a service physically more demanding to provide, a higher recorded burden must not become a verdict that the person deserves less support. Recording resource use and deciding entitlements are separate tasks. The commitment to essential access needs an institutional expression of its own.

\section{What an account could contribute}
A useful account might show that an unreliable pipe repeatedly forces emergency deliveries, or that one area faces avoidable losses. Whether an EBU account actually improves those decisions must be tested. Even a physically accurate calculation cannot itself confer a right, repair exclusion or distribute political power.

\practice{A full reservoir serves an empty clinic tank. List three possible explanations before calling the problem physical scarcity.}{The connecting pump might be broken; the route might not have enough delivery capacity during the relevant interval; or access might be restricted by an institutional or affordability rule. The reservoir's quality and other demands also need checking. The facts given do not choose among these explanations.}

\carry{Scarcity has physical and institutional dimensions, and burdens fall unevenly. These differences matter when we turn to environmental changes that affect the conditions within which people act.}
